\section{Systems and environment}
\label{sec:systems}

\begin{table}[ht]
\centering\footnotesize
\caption{The two systems. Bandwidth and latency were measured by us (Section~\ref{sec:model}).}
\label{tab:systems}
\begin{tabular}{@{}p{3.3cm}p{5.7cm}p{5.4cm}@{}}
\toprule
 & Laptop & Hábrók node (\texttt{regular}) \\
\midrule
CPU & Intel Core i7-12650H & 2 $\times$ AMD EPYC 7763 (Zen 3) \\
Cores & 6 P-cores + 4 E-cores (16 threads), WSL2 shows 10 vCPUs & 2 $\times$ 64 cores, no SMT, 2 NUMA nodes \\
Caches & L1d \SI{48}{KiB} per core, L3 \SI{24}{MiB} & L2 \SI{512}{KiB} per core, L3 16 $\times$ \SI{32}{MiB} (one per CCD) \\
Memory & one DDR4-3200 channel (\SI{25.6}{GB/s} peak), \SI{7.6}{GiB} visible in WSL2 & \SI{503}{GiB} \\
STREAM triad, all / 1 thread & \SI{15.9}{GB/s} / \SI{12.3}{GB/s} & \SI{327}{GB/s} / \SI{34.7}{GB/s} \\
Load latency at \SI{1}{GB} & \SI{376}{ns} & \SI{99}{ns} \\
OS & Ubuntu 24.04.1 on WSL2, Linux 6.18 & AlmaLinux 9.8 \\
Compiler & GCC 13.3, \texttt{-O3 -fopenmp} & GCC 11.5, \texttt{-O3 -fopenmp} \\
Transparent huge pages & \texttt{madvise} (default) & \texttt{always} \\
Tools & perf 6.18 (software events and sampling), gprof & PAPI 7.1 (built by us), pecount, gprof \\
\bottomrule
\end{tabular}
\end{table}

Table~\ref{tab:systems} shows the two systems. All builds also use \texttt{-g
-fno-omit-frame-pointer}, so that perf can record call graphs. This does not change the code
of the loops.

\paragraph{Laptop under WSL2.} WSL2 runs Linux in a Hyper-V virtual machine, and this limits
what we can measure on the laptop:
\begin{itemize}
  \item The hardware counters are not passed through to the virtual machine. All hardware
        events (cycles, instructions, cache and TLB misses) give \emph{not supported}. Only
        software events and timer-based sampling work.
  \item The virtual machine shows 10 vCPUs and does not tell which ones are P-cores and
        which are E-cores, so we cannot choose between them.
  \item The memory addresses are translated twice (nested paging). This is probably one
        reason for the high memory latency of \SI{376}{ns}, and it also adds noise.
\end{itemize}
The laptop has only one memory channel with a peak of \SI{25.6}{GB/s}. STREAM reaches
\SI{15.9}{GB/s} with all threads, which is \SI{62}{\percent} of the peak. In Assignment 6 we
saw on the same laptop that STREAM is fastest with about four threads and gets slower with
more threads. We wanted to boot a native Linux on the laptop to get hardware counters, but we
did not manage this in time, so all hardware counters come from Hábrók.

\paragraph{Hábrók.} The compute nodes have no \texttt{perf}, no \texttt{numactl} and no PAPI
module, and \texttt{perf\_event\_paranoid} is set to 2, so we can only count events in user
space. We therefore built PAPI 7.1 from source in our home directory. We also wrote
\texttt{pecount} (\texttt{pe/bench/pecount.c}, about 200 lines), which counts the events of a
program with \texttt{perf\_event\_open}, like \texttt{perf stat} does. All Hábrók runs used one
full node (exclusive) through Slurm. Still, when several counters were open, they were only
active \SIrange{14}{26}{\percent} of the time. Probably other software on the node also uses
the counters. For exact counts we therefore use PAPI regions with only a few events.
